\section{Discussion}

\subsection{Key Takeaways}

\subsubsection{Takeaway 1: Joint SVD is Efficient}

Sharing the input basis across Q, K, V yields \textbf{{{JOINT\_GAIN}}\%} fewer 
parameters than independent SVD at the same accuracy.

\subsubsection{Takeaway 2: Nested Detection Matters}

The detection order fix restored DistilBERT from 50.92\% collapse to 
{{D\_ACC}}\%.

\subsubsection{Takeaway 3: LoRA Recovery is Essential}

Without LoRA, accuracy drops by {{NOLORA\_DROP}}\%. With LoRA (rank 16), it 
recovers most of the gap.

\subsection{When to Use Compressed MHA}

\subsubsection{Recommended}
\begin{itemize}
\item \textbf{Memory-constrained deployment}
\item \textbf{Fine-tuning pipelines} with LoRA
\item \textbf{Multi-architecture} codebases
\item \textbf{Base for quantization} (compress then quantize)

\end{itemize}
\subsubsection{Not Recommended}
\begin{itemize}
\item \textbf{Latency-critical}: two matmuls vs one
\item \textbf{Extremely small models} ($d < 256$)
\item \textbf{Extremely large rank needed}: $r > d/2$

\end{itemize}
\subsection{Decision Flowchart}

\begin{lstlisting}
Want to compress attention projections? --- Yes --- Use Compressed MHA
  | No
   
Need sequence-level compression? --- Yes --- Use Linformer/LEKAv3
  | No
   
Keep standard attention
\end{lstlisting}

\subsection{Cost-Benefit}

For DistilBERT with rank 64:

\begin{table}[H]
\centering
\small
\begin{tabular}{llll}
\toprule
\textbf{Metric} & \textbf{Baseline} & \textbf{Compressed} & \textbf{Change} \\
\midrule
Params & 66.95M & {{D\_PARAMS}}M & \textbf{{{D\_PARAM\_PCT}}\%} \\
Attention params & 14.2M & {{D\_ATTN\_PARAMS}}M & \textbf{{{D\_ATTN\_RATIO}}$\times$} \\
Peak VRAM & 183 MB & {{D\_VRAM}} MB & {{D\_VRAM\_PCT}}\% \\
Accuracy & 90.71\% & {{D\_ACC}}\% & {{D\_DELTA}}\% \\
\bottomrule
\end{tabular}
\end{table}

\subsection{Comparison with Prior Compression Methods}

\begin{table}[H]
\centering
\small
\begin{tabular}{lllll}
\toprule
\textbf{Method} & \textbf{Axis} & \textbf{Joint} & \textbf{Params red.} & \textbf{Accuracy} \\
\midrule
Linformer & Sequence & No & No & {{LIN\_ACC}} \\
A3 (SVD) & Feature & No & Yes & {{A3\_ACC}} \\
AASVD & Feature & No & Yes & {{AASVD\_ACC}} \\
\textbf{Compressed MHA} & \textbf{Feature (joint)} & \textbf{Yes} & \textbf{Yes} & \textbf{{{D\_ACC}}} \\
\bottomrule
\end{tabular}
\end{table}

\subsection{Deployment Considerations}

\subsubsection{Memory}

Compression ratio {{D\_RATIO}}$\times$ reduces both disk and VRAM footprint.

\subsubsection{Latency}

Two matmuls per projection vs one for baseline. On modern GPUs, the smaller 
matrices are often faster due to reduced memory bandwidth.

\subsubsection{Batching}

Compression scales well with batch size -- the smaller weight matrices fit 
better in cache.

\subsection{Open Questions}

\begin{enumerate}
\item \textbf{Optimal rank allocation:} Per-layer rank schedules?
\item \textbf{Joint with output:} Can we joint-factorize Q/K/V/O?
\item \textbf{Cross-layer SVD:} Share factors across layers?
\item \textbf{Interaction with quantization:} Combined INT8 + SVD?

\end{enumerate}
\subsection{Broader Impact}

\subsubsection{Positive}
\begin{itemize}
\item Reduces model size for edge deployment
\item Complements other compression methods
\item Multi-architecture support aids adoption

\end{itemize}
\subsubsection{Negative}
\begin{itemize}
\item Requires fine-tuning for recovery
\item Adds complexity vs simpler pruning
\end{itemize}
\begin{lstlisting}

---
\end{lstlisting}